\section{Opening the slab}\label{sec:opening}

The local estimates and cell paths now define a change from the closed
Gibbs law to an open path measure. We specify the target measure, the
accepted choices, and the interacting proposal in that order. Keeping
its normalization exact will let the two copies of the change cancel
in the next section.

Let \(q_0=\langle u_0,\gamma^{(1)}u_0\rangle/N\). Starting from
\eqref{eq:14}, average a slot \(i_0\) with factor \(1/N\) and replace
its inner end by \(E_{i_0}=z\), integrated with \(dz/V\). All other
inner ends remain \(Y_i\), and every outer endpoint \(Y_i\) is still
integrated. The resulting finite measure has total mass \(q_0\): it
inserts \(N^{-1}\sum_i(|u_0\rangle\langle u_0|)_i\) between the two
kernels in the trace. Add a dummy point \(x\), uniform in the cell of
\(Y_{i_0}\), and retain \(i_0\) as a coordinate. Denote this augmented
measure by \(\mu\); the dummy coordinate does not change its mass.

\subsection{The base law and acceptance}
Our base law is the closed probability \eqref{eq:14}, denoted by the
configuration $C$, together with independent uniform points $x,z$ of
law $dx\,dz/V^2$. In the slab notation, configurations are tuples of
paths rather than density operators; their starting tuples carry
superscripts, such as $S^C$, and a subscript restricts their paths to
a set of slots.

Set $a=\operatorname{cell}(x)$ and $b=\operatorname{cell}(z)$.
The first index of $\mu_{ab}$ in Section~\ref{sec:defects} specifies
the new inner endpoint's cell, and the second specifies the outer
endpoint's cell. Thus the two orientations used below are
\[
\begin{array}{c|c|c|c}
 &\text{opened slot}&\text{outer endpoint}&\text{new inner endpoint}\\ \hline
 \text{forward insertion }\mu_{ba}&i_0&Y_{i_0}\in a&z\in b\\
 \text{reverse insertion }\mu_{ab}&e&Y_e\in b&x\in a.
\end{array}
\]
The reverse slot $e$ will be the last label on the selected path.
Only valid ordered pairs, meaning the admissible pairs of
Section~\ref{sec:skeletons}, are used. For such a pair draw a skeleton
and apply the deterministic adaptation to the unusable cells of $C$.
Require geometric acceptance and the two tests
$\mathcal A^+(x;C),\mathcal A^+(z;C)\le C_0$.
Denote these combined conditions by $\operatorname{acc}(C,\gamma)$.
For an invalid pair the acceptance indicator and all weights below
are zero. The base average acceptance mass is at least $\upsilon/2$:
the valid-pair fraction is at least $\upsilon$, and the geometric
rejection and dummy-score probabilities can each be made smaller
than a fixed fraction of it by the choices already specified.

\subsection{Selecting one eligible slot in each path cell}
On acceptance let $P(C,\gamma)$ be the ordered simple star path.
For a cell $Q$, write
\[
 g_Q(C)=\sum_{i:Y_i\in Q}\mathbf1_{\{i\text{ eligible in }C\}}.
\]
Independently in each path cell choose a uniform eligible slot.
This gives an ordered tuple $I=(i_0,\ldots,i_k)$, one label per
cell. For any candidate tuple define, on acceptance,
\[
\begin{split}
 h_\gamma(C,I)={}&
 \mathbf1_{\{(\operatorname{cell}(Y_{i_0}),\ldots,
                      \operatorname{cell}(Y_{i_k}))=P(C,\gamma)\}}\\
 &\quad\cdot\prod_{j=0}^k
       \frac{\mathbf1_{\{i_j\text{ eligible in }C\}}}
            {g_{\operatorname{cell}(Y_{i_j})}(C)}.
\end{split}
\]
Set the weight to zero off acceptance, and interpret a factor with
zero eligible count as zero. The cells of $P$ are distinct and
usable, so on acceptance the displayed factors are defined and
$\sum_Ih_\gamma(C,I)=1$. Equivalently, the sum ranges over candidate
simple star paths from $a$ to $b$ and label choices in their outer
$Y$-cells. This range depends on $a,b,Y$, whereas the weight retains
all dependence on the actual paths of $C$.

\subsection{The conditional interacting proposal}
For a tuple $I$, use cyclic predecessor and successor, acting as the
identity off the tuple, and let \((p_I S)_i=S_{\operatorname{pred}_I(i)}\).
Transform the starts to \(S^O=p_I S^C\). The inner ends \(E^O\) agree
with $Y$ except at slot \(i_0\), whose end is replaced by \(z\);
paths off $I$ are kept. Resample paths in \(I\) from their conditional
interacting distribution restricted by \(\mathrm{cut}_R\).

More precisely, for a state \(X\) with fixed endpoints and complement
paths, let \(W_G^X\) use the free bridge measures on \(G\) and the
exponential of minus the interactions incident to \(G\), each pair once.
Write \(\mathcal Z_G(X)=\int\mathbf 1_{\mathrm{cut}_R(G)} W_G^X\).
The proposal on \(I\) uses this cut measure in \(O\), divided by
\(\mathcal Z_I(O)\). The new displacements are at most \((D+10)s\)
by eligibility and star adjacency: the start from the last cell goes to
the end \(z\). Thus the normalization is positive.

\subsection{The density and its second moment}

Averaging a path-dependent likelihood before estimating its second
moment also appears in random-trail detection and perturbed-lattice
reassignment~\cite{AriasCastroCandesHelgasonZeitouni2008,
BergerPeres2013,PeresSly2014}. In the present setting the conditional
interacting measures and both proposal normalizers must be retained
explicitly.

Write \(\nu\) for the resulting subprobability measure on the open
variables. Its mass is the base acceptance probability, hence at least
\(\upsilon/2\). Its density relative to \(\mu\) is
\[
 f = m\,\mathbb E_\gamma\sum_{I:\,\mathrm{first}(I)=i_0}
       \frac{\mathbf 1_{\mathrm{cut}_R(O_I)}}{\mathcal Z_I(O)}
              \int W_I^{\widehat C}\, h_\gamma(\widehat C,I),                             \tag{26}\label{eq:26}
\]
where the hypothetical closed state \(\widehat C\) has
\(S^{\widehat C}=p_I^{-1}S^O\), \(E^{\widehat C}=Y\), and paths off
\(I\) as in \(O\). A term with zero old weight is set to zero even
if its normalizer vanishes; positive old weight gives the displacement
condition and a positive normalizer.

To verify the density, change start variables for each candidate tuple.
The Jacobian is one and the outer kernel is invariant under this
permutation. Write \(W_{\bar I}^{\rm int}\) for the complement bridge
measures with only their internal interactions. This factor agrees
before and after the change. Integrating the old paths on \(I\) leaves
\[
 W_{\bar I}^{\rm int}
   \left[\int W_I^{\widehat C} h_\gamma(\widehat C,I)\right]
     \mathbf 1_{\mathrm{cut}_R(O_I)}W_I^O/\mathcal Z_I(O)
\]
(all incident weights and normalizers use the displayed fixed complement).
The corresponding factor in the open measure is
\(W_{\bar I}^{\rm int}W_I^O\). The remaining normalization comes from
the label and dummy coordinate:
\[
 \frac{\text{base density of }x}
      {\text{target index factor}\times\text{target density of }x}
 =\frac{V^{-1}}{N^{-1}s^{-3}}=\frac{N}{M}=m.
\]
Here both measures restrict $x$ to the cell containing $Y_{i_0}$,
and both use $dz/V$. This proves \eqref{eq:26}, equivalently by
integration against nonnegative test functions. Skeleton expectations
are for the given cell pair, with zero weight for invalid pairs.

Cauchy--Schwarz gives
\((\upsilon/2)^2\le q_0\int f^2\dd\mu\). To compute the second
moment, use \(\int f^2\dd\mu=\int f\dd\nu\): evaluate
\eqref{eq:26} at a transformation from the base with independent
skeleton \(\gamma'\) and tuple \(I'\). The resulting integral is
\[
 \begin{split}
 &m\,\mathbb E_{x,z,\gamma',\gamma}\, Z_N^{-1}\int
 F(Y,S^C)\,dY\,dS^C
 \sum_{I',I:\,\mathrm{first}(I')=\mathrm{first}(I)}\int
   W(C) h_{\gamma'}(C,I')\\
 &\hspace{30mm}\cdot
     \frac{\mathbf 1_{\mathrm{cut}_R(O_{I'})}W_{I'}^O}
          {\mathcal Z_{I'}(O)}
     \frac{\mathbf 1_{\mathrm{cut}_R(O_I)}}{\mathcal Z_I(O)}
          W_I^{\widehat C} h_\gamma(\widehat C,I).
 \end{split}                                                          \tag{27}\label{eq:27}
\]
The path integrations are over all closed paths, then proposed paths
on \(I'\), and finally hypothetical old paths on \(I\). Here \(W(C)\)
is the full slab weight, \(S^O=p_{I'}S^C\), and $x,z$ have their base
uniform law. Both tuples must have the same first slot, which is summed
over. All integrands are nonnegative, so Tonelli permits the changes in
integration order used next. Section~\ref{sec:comparison} bounds
\eqref{eq:27} uniformly.
